% generated by docs/gen_function_reference.py -- do not edit by hand
\section{Complete function reference}\label{app:functions}
Every function and class in the project (tests excluded), grouped by package and module, generated from the source code so that nothing is missed. Methods are listed as \texttt{Class.method}. In total: 284 entries in 70 modules.

\subsection*{\texttt{data/}}
\begin{longtable}{@{}>{\raggedright\arraybackslash}p{5.6cm}>{\raggedright\arraybackslash}p{10.6cm}@{}}
\multicolumn{2}{@{}l}{\small\bfseries\texttt{data/generate\_diffusion.py}}\\[1pt]
\quad\fn{load_config} & \small Read config.yaml into a dict. \\
\quad\fn{initial_condition} & \small The starting profile u(x,0): one sine (degenerate), two sines (default) or a Gaussian bump. \\
\quad\fn{solve_diffusion_fd} & \small Explicit finite-difference solver for u\_t = D u\_xx with zero boundary values; produces the ground truth. \\
\quad\fn{build_full_field} & \small Build the x and t grids from config and run the solver: the clean 800x61 field. \\
\quad\fn{flatten_field} & \small Turn the 2-D field into matching 1-D arrays (x, t, u), one row per grid point. \\
\quad\fn{make_sparse_and_noisy} & \small Draw the 10\% noise-free subsample and the 400 noisy observations (5\% noise). \\
\quad\fn{main} & \small Phase 1 entry point: generate and save full.npz, sparse.npz and noisy.npz. \\
\addlinespace[4pt]
\multicolumn{2}{@{}l}{\small\bfseries\texttt{data/problem\_description.py}}\\[1pt]
\quad\fn{build_problem_description} & \small Compute statistics of the noisy observations (amplitude change, edges, centroid, smoothness) and write them as plain text, never naming a mechanism. \\
\addlinespace[4pt]
\end{longtable}
\subsection*{\texttt{equation\_discovery/}}
\begin{longtable}{@{}>{\raggedright\arraybackslash}p{5.6cm}>{\raggedright\arraybackslash}p{10.6cm}@{}}
\multicolumn{2}{@{}l}{\small\bfseries\texttt{equation\_discovery/derivatives.py}}\\[1pt]
\quad\fn{reconstruct_grid} & \small Interpolate scattered (x, t, u) samples onto a regular grid (scipy griddata, nearest-neighbour at edges). \\
\quad\fn{smooth_grid} & \small Gaussian smoothing of a reconstructed grid to reduce noise before differentiation. \\
\quad\fn{finite_diff_derivatives} & \small Central finite differences on a grid: u\_t, u\_x, u\_xx, u\_xxx, with edge cropping. \\
\addlinespace[4pt]
\multicolumn{2}{@{}l}{\small\bfseries\texttt{equation\_discovery/discovery\_engine.py}}\\[1pt]
\quad\fn{_fields_from_grid} & \small Compute derivative fields on a grid and flatten them for regression. \\
\quad\fn{discover_from_grid} & \small Full discovery on gridded data: derivatives, library, sparse regression. \\
\quad\fn{discover_from_scatter} & \small Full discovery on scattered data: reconstruct, optionally smooth, then discover\_from\_grid. \\
\quad\fn{equation_to_string} & \small Format a discovery result as a readable equation string. \\
\quad\fn{evaluate_against_ground_truth} & \small Compare a discovered equation with the hidden truth (grading only, after discovery). \\
\quad\fn{evaluate_competing_hypotheses} & \small Fit each named single-term hypothesis separately and report R\^{}2 and coefficients. \\
\addlinespace[4pt]
\multicolumn{2}{@{}l}{\small\bfseries\texttt{equation\_discovery/library.py}}\\[1pt]
\quad\fn{term_label} & \small Human-readable label for a library term, e.g. ('u','u\_x') -> 'u*u\_x'. \\
\quad\fn{build_library} & \small Evaluate every candidate term at every sample: the matrix Theta. \\
\quad\fn{parse_term_specs} & \small Convert the config's term lists into term tuples. \\
\addlinespace[4pt]
\multicolumn{2}{@{}l}{\small\bfseries\texttt{equation\_discovery/sindy.py}}\\[1pt]
\quad\fn{stlsq} & \small Sequentially thresholded least squares: fit, zero small coefficients, refit, repeat (columns normalised). \\
\quad\fn{run_sindy} & \small Fit u\_t = Theta xi with STLSQ (pysindy if installed, else stlsq) and return a structured result. \\
\quad\fn{fit_restricted} & \small Ordinary least squares on a fixed set of terms (used for competing hypotheses). \\
\addlinespace[4pt]
\end{longtable}
\subsection*{\texttt{pinn/}}
\begin{longtable}{@{}>{\raggedright\arraybackslash}p{5.6cm}>{\raggedright\arraybackslash}p{10.6cm}@{}}
\multicolumn{2}{@{}l}{\small\bfseries\texttt{pinn/data\_utils.py}}\\[1pt]
\quad\fn{_to_tensor} & \small Convert a numpy array to a float32 column tensor. \\
\quad\fn{load_observational_data} & \small Load noisy.npz or sparse.npz, returning only x, t, u (the only data a PINN may train on). \\
\quad\fn{load_clean_reference} & \small Load the clean field for evaluation only, never for training. \\
\quad\fn{make_collocation_points} & \small Random (x, t) points where the physics residual is penalised. \\
\quad\fn{make_bc_points} & \small Random points on both boundaries with target u = 0. \\
\quad\fn{make_ic_points} & \small Random points at t = 0 with target u = initial condition. \\
\quad\fn{data_dict_to_tensors} & \small Convert an (x, t, u) dict into tensors. \\
\quad\fn{filter_by_time} & \small Keep only observations with t below a cutoff (train/OOD split). \\
\addlinespace[4pt]
\multicolumn{2}{@{}l}{\small\bfseries\texttt{pinn/derivatives.py}}\\[1pt]
\quad\fn{compute_derivatives} & \small Exact u, u\_t, u\_x, u\_xx (and u\_xxx) of the network output via torch.autograd. \\
\addlinespace[4pt]
\multicolumn{2}{@{}l}{\small\bfseries\texttt{pinn/forward.py}}\\[1pt]
\quad\fn{run_forward_pinn} & \small Train a PINN with physics parameters fixed (known D). \\
\addlinespace[4pt]
\multicolumn{2}{@{}l}{\small\bfseries\texttt{pinn/inverse.py}}\\[1pt]
\quad\fn{run_inverse_pinn} & \small Train a PINN with physics parameters trainable (estimate D). \\
\addlinespace[4pt]
\multicolumn{2}{@{}l}{\small\bfseries\texttt{pinn/losses.py}}\\[1pt]
\quad\fn{mse} & \small Mean squared error helper. \\
\quad\fn{data_loss} & \small Mismatch between network and observations. \\
\quad\fn{physics_loss} & \small Mean squared equation residual at collocation points. \\
\quad\fn{bc_loss} & \small Mismatch at the boundaries. \\
\quad\fn{ic_loss} & \small Mismatch with the initial condition. \\
\quad\fn{total_loss} & \small Weighted sum of the four losses. \\
\addlinespace[4pt]
\multicolumn{2}{@{}l}{\small\bfseries\texttt{pinn/network.py}}\\[1pt]
\quad\fn{PINN} & \small The one reusable network: (x, t) -> u through four tanh layers of 64 units. \\
\quad\fn{PINN._init_weights} & \small Xavier initialisation of the weights. \\
\quad\fn{PINN.forward} & \small Evaluate u at given (x, t). \\
\addlinespace[4pt]
\multicolumn{2}{@{}l}{\small\bfseries\texttt{pinn/residuals.py}}\\[1pt]
\quad\fn{diffusion_residual} & \small Residual u\_t - D u\_xx. \\
\quad\fn{advection_residual} & \small Residual u\_t - c u\_x. \\
\quad\fn{reaction_diffusion_residual} & \small Residual u\_t - D u\_xx + k u. \\
\quad\fn{make_params} & \small Wrap initial parameter values as tensors, trainable or fixed. \\
\addlinespace[4pt]
\multicolumn{2}{@{}l}{\small\bfseries\texttt{pinn/trainer.py}}\\[1pt]
\quad\fn{PINNTrainer} & \small Holds network, residual function, parameters and loss weights; runs training. \\
\quad\fn{PINNTrainer.trainable_parameters} & \small Network weights plus any trainable physics parameters. \\
\quad\fn{PINNTrainer.compute_losses} & \small One evaluation of all loss terms (derivatives, residual, data, BC, IC). \\
\quad\fn{PINNTrainer.train} & \small Adam steps then an L-BFGS polish, recording loss history. \\
\quad\fn{PINNTrainer._record} & \small Store one row of loss history and current parameter values. \\
\quad\fn{build_and_train} & \small Shared entry point: build network and parameters, train, return trainer and history. \\
\addlinespace[4pt]
\end{longtable}
\subsection*{\texttt{validation/}}
\begin{longtable}{@{}>{\raggedright\arraybackslash}p{5.6cm}>{\raggedright\arraybackslash}p{10.6cm}@{}}
\multicolumn{2}{@{}l}{\small\bfseries\texttt{validation/conservation.py}}\\[1pt]
\quad\fn{conservation_violation} & \small Global balance check: does the model's total amount change exactly as its own law says (dM/dt = integral of f)? \\
\addlinespace[4pt]
\multicolumn{2}{@{}l}{\small\bfseries\texttt{validation/forward\_solver.py}}\\[1pt]
\quad\fn{rhs_function} & \small Compile any candidate equation into a numerical right-hand side u\_t = f(u, u\_x, u\_xx). \\
\quad\fn{_effective_coefficients} & \small Estimate diffusion/advection/reaction strengths to choose a stable time step. \\
\quad\fn{solve_forward} & \small Solve a candidate law forward in time from an initial condition (Euler or RK4). \\
\addlinespace[4pt]
\multicolumn{2}{@{}l}{\small\bfseries\texttt{validation/model\_selection.py}}\\[1pt]
\quad\fn{parameter_contributions} & \small Term necessity: share of RMS(u\_t) contributed by each parameter's terms on the trained PINN. \\
\quad\fn{reduced_signature} & \small Signature of the equation with inactive terms removed (does it reduce to another candidate?). \\
\quad\fn{bic} & \small Bayesian information criterion n ln(MSE) + k ln n. \\
\quad\fn{select_model} & \small The selection rule: gate, fit-equivalence, optional reduction, complexity, BIC. \\
\addlinespace[4pt]
\multicolumn{2}{@{}l}{\small\bfseries\texttt{validation/ood.py}}\\[1pt]
\quad\fn{split_by_time} & \small Split points into t <= cutoff and t > cutoff. \\
\quad\fn{prediction_error} & \small RMSE and MAE of a model against reference values. \\
\addlinespace[4pt]
\multicolumn{2}{@{}l}{\small\bfseries\texttt{validation/physics\_validation.py}}\\[1pt]
\quad\fn{evaluate_physics_residual} & \small Mean squared residual of a trained PINN on a fresh grid (not the training points). \\
\addlinespace[4pt]
\multicolumn{2}{@{}l}{\small\bfseries\texttt{validation/scorer.py}}\\[1pt]
\quad\fn{build_validation_result} & \small Turn measured losses into supported/rejected using fixed thresholds (Phase 3). \\
\addlinespace[4pt]
\multicolumn{2}{@{}l}{\small\bfseries\texttt{validation/stability.py}}\\[1pt]
\quad\fn{candidate_monomials} & \small The distinct field terms on a candidate's right-hand side (e.g. u\_xx, u, u\^{}2). \\
\quad\fn{bootstrap_stability} & \small Refit those terms by least squares on 20 random 80\% subsets; report mean, spread, sign stability. \\
\addlinespace[4pt]
\end{longtable}
\subsection*{\texttt{llm/}}
\begin{longtable}{@{}>{\raggedright\arraybackslash}p{5.6cm}>{\raggedright\arraybackslash}p{10.6cm}@{}}
\multicolumn{2}{@{}l}{\small\bfseries\texttt{llm/\_\_init\_\_.py}}\\[1pt]
\quad\fn{get_provider} & \small Choose the LLM provider from config and environment (Gemini if a key is set, else the mock). \\
\addlinespace[4pt]
\multicolumn{2}{@{}l}{\small\bfseries\texttt{llm/base.py}}\\[1pt]
\quad\fn{LLMProvider} & \small Abstract interface every LLM provider implements. \\
\quad\fn{LLMProvider.generate} & \small Return the raw text completion for a prompt. \\
\quad\fn{LLMProvider.name} & \small The provider's class name, recorded in results. \\
\addlinespace[4pt]
\multicolumn{2}{@{}l}{\small\bfseries\texttt{llm/gemini\_provider.py}}\\[1pt]
\quad\fn{_repair_json_escapes} & \small If the response is not valid JSON, double stray backslashes (LaTeX like \textbackslash{}partial). \\
\quad\fn{GeminiProvider} & \small Real provider for Google Gemini; the key comes only from GEMINI\_API\_KEY. \\
\quad\fn{GeminiProvider._get_client} & \small Lazily import google-genai and create the client; clear error if no key. \\
\quad\fn{GeminiProvider.generate} & \small Call Gemini in JSON mode, retrying overload errors (429/500/503) with back-off. \\
\addlinespace[4pt]
\multicolumn{2}{@{}l}{\small\bfseries\texttt{llm/mock\_provider.py}}\\[1pt]
\quad\fn{MockLLMProvider} & \small Deterministic offline provider returning a canned set of four hypotheses. \\
\quad\fn{MockLLMProvider.generate} & \small Return the same canned JSON every time, whatever the prompt. \\
\addlinespace[4pt]
\multicolumn{2}{@{}l}{\small\bfseries\texttt{llm/openai\_provider.py}}\\[1pt]
\quad\fn{OpenAIProvider} & \small Legacy OpenAI provider (no longer selectable by get\_provider). \\
\quad\fn{OpenAIProvider._get_client} & \small Lazily create the OpenAI client. \\
\quad\fn{OpenAIProvider.generate} & \small Call OpenAI chat completions. \\
\addlinespace[4pt]
\multicolumn{2}{@{}l}{\small\bfseries\texttt{llm/prompts.py}}\\[1pt]
\quad\fn{build_hypothesis_prompt_with_context} & \small Mode B prompt: observations plus retrieved evidence, mechanisms, equations and parameters. \\
\quad\fn{build_hypothesis_prompt} & \small Mode A prompt: observations only, with the JSON schema the answer must follow. \\
\addlinespace[4pt]
\multicolumn{2}{@{}l}{\small\bfseries\texttt{llm/schemas.py}}\\[1pt]
\quad\fn{ParameterSpec} & \small One parameter: name, trainable flag, initial value, units, optional bounds. \\
\quad\fn{Hypothesis} & \small The structured hypothesis every LLM answer must fit (pydantic model). \\
\quad\fn{Hypothesis._confidence_in_range} & \small Validator: confidence must lie in [0, 1]. \\
\quad\fn{Hypothesis._equation_has_equals} & \small Validator: the equation must contain '='. \\
\quad\fn{HypothesisSet} & \small A validated collection of hypotheses for one problem. \\
\addlinespace[4pt]
\end{longtable}
\subsection*{\texttt{hypotheses/}}
\begin{longtable}{@{}>{\raggedright\arraybackslash}p{5.6cm}>{\raggedright\arraybackslash}p{10.6cm}@{}}
\multicolumn{2}{@{}l}{\small\bfseries\texttt{hypotheses/adapter.py}}\\[1pt]
\quad\fn{hypothesis_to_pinn_config} & \small Compile a hypothesis into exactly what build\_and\_train needs (residual function, parameters). \\
\addlinespace[4pt]
\multicolumn{2}{@{}l}{\small\bfseries\texttt{hypotheses/complexity.py}}\\[1pt]
\quad\fn{equation_complexity} & \small Complexity score: terms + free parameters + derivative order + 2 x nonlinear + 0.1 x operations. \\
\addlinespace[4pt]
\multicolumn{2}{@{}l}{\small\bfseries\texttt{hypotheses/constraints.py}}\\[1pt]
\quad\fn{resolve_bounds} & \small Combine a hypothesis's declared bounds with the config fallbacks (D > 0, K > 0). \\
\quad\fn{precheck_physics} & \small Before training: reject non-evolution equations; warn on over-determined or uncoupled boundary setups. \\
\quad\fn{check_fitted_bounds} & \small After training: flag fitted parameters that violate physical bounds. \\
\addlinespace[4pt]
\multicolumn{2}{@{}l}{\small\bfseries\texttt{hypotheses/equation\_compiler.py}}\\[1pt]
\quad\fn{EquationParseError} & \small Raised when an equation string cannot be parsed. \\
\quad\fn{parse_equation} & \small Parse 'lhs = rhs' with SymPy into a residual; detect fields and parameters used. \\
\quad\fn{compile_residual} & \small Turn the residual into a torch function residual\_fn(derivs, params) for the trainer. \\
\quad\fn{equation_str_to_sympy_str} & \small Normalised human-readable form of an equation. \\
\addlinespace[4pt]
\multicolumn{2}{@{}l}{\small\bfseries\texttt{hypotheses/pipeline.py}}\\[1pt]
\quad\fn{_extract_json_array} & \small Strip markdown code fences and parse the LLM's JSON array. \\
\quad\fn{generate_hypotheses} & \small Prompt, call the provider, parse, schema-check, validate, deduplicate and rank. \\
\addlinespace[4pt]
\multicolumn{2}{@{}l}{\small\bfseries\texttt{hypotheses/ranking.py}}\\[1pt]
\quad\fn{_completeness_score} & \small Fraction of descriptive fields the LLM filled in. \\
\quad\fn{_simplicity_score} & \small 1 / (1 + number of parameters). \\
\quad\fn{rank_hypotheses} & \small Composite score that sets investigation order only, never the verdict. \\
\addlinespace[4pt]
\multicolumn{2}{@{}l}{\small\bfseries\texttt{hypotheses/search.py}}\\[1pt]
\quad\fn{kg_template_candidates} & \small Turn equations documented in the knowledge graph into candidate hypotheses. \\
\quad\fn{build_candidate_pool} & \small Merge LLM and graph candidates and remove structural duplicates. \\
\quad\fn{structural_and_physics_filter} & \small Run deterministic validation and physics pre-checks; keep the survivors. \\
\quad\fn{make_splits} & \small Train / random holdout / OOD (t > cutoff) splits of the observations. \\
\quad\fn{_rmse} & \small RMSE of a network on a split. \\
\quad\fn{validate_candidate} & \small Inverse PINN on the training split, then gate checks, bounds, term necessity, complexity, BIC. \\
\quad\fn{load_candidate_model} & \small Reload a saved candidate network. \\
\quad\fn{train_data_only_mlp} & \small Baseline: the same network trained on data only, without physics. \\
\quad\fn{save_json} & \small Write a result dict to JSON, creating folders as needed. \\
\addlinespace[4pt]
\multicolumn{2}{@{}l}{\small\bfseries\texttt{hypotheses/validation.py}}\\[1pt]
\quad\fn{validate_hypothesis} & \small Deterministic checks: syntax, variables, parameters, derivatives, units. \\
\quad\fn{canonical_signature} & \small Parameter- and sign-invariant structure of an equation, for duplicate detection. \\
\quad\fn{deduplicate} & \small Remove hypotheses with identical canonical signatures. \\
\addlinespace[4pt]
\end{longtable}
\subsection*{\texttt{rag/}}
\begin{longtable}{@{}>{\raggedright\arraybackslash}p{5.6cm}>{\raggedright\arraybackslash}p{10.6cm}@{}}
\multicolumn{2}{@{}l}{\small\bfseries\texttt{rag/chunker.py}}\\[1pt]
\quad\fn{chunk_document} & \small Split a document into chunks at its section headers. \\
\addlinespace[4pt]
\multicolumn{2}{@{}l}{\small\bfseries\texttt{rag/context.py}}\\[1pt]
\quad\fn{ScientificContext} & \small The unified evidence object: retrieved evidence, mechanisms, equations, parameters, sources. \\
\quad\fn{build_scientific_context} & \small Retrieve evidence for the description, then query the graph for every mechanism it mentions. \\
\addlinespace[4pt]
\multicolumn{2}{@{}l}{\small\bfseries\texttt{rag/corpus\_loader.py}}\\[1pt]
\quad\fn{_parse_frontmatter} & \small Split a Markdown file into its metadata header and body. \\
\quad\fn{load_document} & \small Load one corpus file as a Document. \\
\quad\fn{load_corpus} & \small Load every .md file in rag/corpus/. \\
\addlinespace[4pt]
\multicolumn{2}{@{}l}{\small\bfseries\texttt{rag/documents.py}}\\[1pt]
\quad\fn{DocumentMetadata} & \small A document's id, title, source, domain and declared mechanism, equations and parameters. \\
\quad\fn{Document} & \small Metadata plus text. \\
\quad\fn{Chunk} & \small A piece of a document with its section and id. \\
\addlinespace[4pt]
\multicolumn{2}{@{}l}{\small\bfseries\texttt{rag/embeddings.py}}\\[1pt]
\quad\fn{EmbeddingProvider} & \small Interface for turning text into vectors. \\
\quad\fn{EmbeddingProvider.embed} & \small Embed one text. \\
\quad\fn{EmbeddingProvider.embed_many} & \small Embed many texts. \\
\quad\fn{EmbeddingProvider.name} & \small Provider name. \\
\quad\fn{LocalEmbeddingProvider} & \small Deterministic offline embedding: words hashed into 256 slots, normalised. \\
\quad\fn{LocalEmbeddingProvider._token_index} & \small Hash a word to a slot (MD5 modulo 256). \\
\quad\fn{LocalEmbeddingProvider.embed} & \small Count words per slot and normalise to unit length. \\
\quad\fn{APIEmbeddingProvider} & \small Optional real embedding API provider (unused by default). \\
\quad\fn{APIEmbeddingProvider._get_client} & \small Lazily create its API client. \\
\quad\fn{APIEmbeddingProvider.embed} & \small Embed via the API. \\
\addlinespace[4pt]
\multicolumn{2}{@{}l}{\small\bfseries\texttt{rag/evidence.py}}\\[1pt]
\quad\fn{RetrievalResult} & \small One retrieved chunk with its similarity score. \\
\quad\fn{Evidence} & \small A retrieved claim with provenance; its confidence is retrieval confidence only. \\
\quad\fn{evidence_from_retrieval} & \small Convert a retrieval result into an Evidence object. \\
\addlinespace[4pt]
\multicolumn{2}{@{}l}{\small\bfseries\texttt{rag/metadata.py}}\\[1pt]
\quad\fn{matches_filters} & \small Check a chunk's metadata against filter values. \\
\addlinespace[4pt]
\multicolumn{2}{@{}l}{\small\bfseries\texttt{rag/pipeline.py}}\\[1pt]
\quad\fn{RAGPipeline} & \small Ties chunker, embedder and vector store together. \\
\quad\fn{RAGPipeline.ingest} & \small Chunk, embed and store every document. \\
\quad\fn{RAGPipeline.get_retriever} & \small Return a Retriever over the stored chunks. \\
\addlinespace[4pt]
\multicolumn{2}{@{}l}{\small\bfseries\texttt{rag/retriever.py}}\\[1pt]
\quad\fn{Retriever} & \small Search interface over the vector store. \\
\quad\fn{Retriever.retrieve} & \small Embed the query, take the most similar chunks, apply filters, return Evidence. \\
\addlinespace[4pt]
\multicolumn{2}{@{}l}{\small\bfseries\texttt{rag/vector\_store.py}}\\[1pt]
\quad\fn{VectorStore} & \small Interface: add, search, get. \\
\quad\fn{VectorStore.add} & \small Store a vector with its payload. \\
\quad\fn{VectorStore.search} & \small Return the top-k most similar vectors. \\
\quad\fn{VectorStore.get} & \small Fetch one stored item. \\
\quad\fn{InMemoryVectorStore} & \small Brute-force cosine-similarity store in memory. \\
\quad\fn{InMemoryVectorStore.add} & \small Store a vector with its payload. \\
\quad\fn{InMemoryVectorStore.search} & \small Cosine similarity against every stored vector; return the top-k. \\
\quad\fn{InMemoryVectorStore.get} & \small Fetch one stored item. \\
\quad\fn{InMemoryVectorStore.all_ids} & \small List stored ids. \\
\quad\fn{PersistentVectorStore} & \small The in-memory store plus JSON save/load. \\
\quad\fn{PersistentVectorStore.save} & \small Write the store to JSON. \\
\quad\fn{PersistentVectorStore.load} & \small Read the store from JSON. \\
\addlinespace[4pt]
\end{longtable}
\subsection*{\texttt{knowledge\_graph/}}
\begin{longtable}{@{}>{\raggedright\arraybackslash}p{5.6cm}>{\raggedright\arraybackslash}p{10.6cm}@{}}
\multicolumn{2}{@{}l}{\small\bfseries\texttt{knowledge\_graph/builder.py}}\\[1pt]
\quad\fn{_slug} & \small Make a stable id fragment from a name. \\
\quad\fn{_find_chunk_id} & \small Find the chunk that states a fact (provenance). \\
\quad\fn{build_graph_from_documents} & \small Create Document, Mechanism, Equation, Variable and Parameter entities and sourced relationships. \\
\addlinespace[4pt]
\multicolumn{2}{@{}l}{\small\bfseries\texttt{knowledge\_graph/graph.py}}\\[1pt]
\quad\fn{_relationship_id} & \small Stable hash of (subject, predicate, object), so repeated facts merge. \\
\quad\fn{KnowledgeGraph} & \small The graph container (backed by networkx). \\
\quad\fn{KnowledgeGraph.add_entity} & \small Add an entity, merging if its id exists. \\
\quad\fn{KnowledgeGraph.get_entity} & \small Look up an entity by id. \\
\quad\fn{KnowledgeGraph.entities_by_type} & \small All entities of one type. \\
\quad\fn{KnowledgeGraph.add_relationship} & \small Add a typed, sourced edge, merging sources for duplicates. \\
\quad\fn{KnowledgeGraph.neighbors} & \small Entities connected to a given entity. \\
\quad\fn{KnowledgeGraph.relationships_for} & \small Edges touching a given entity. \\
\quad\fn{KnowledgeGraph.summary} & \small Counts of entities and relationships by type. \\
\addlinespace[4pt]
\multicolumn{2}{@{}l}{\small\bfseries\texttt{knowledge\_graph/query.py}}\\[1pt]
\quad\fn{find_entity_by_name} & \small Find an entity by its name. \\
\quad\fn{equations_for_mechanism} & \small Documented equations that govern a mechanism. \\
\quad\fn{variables_for_mechanism} & \small Variables in those equations. \\
\quad\fn{parameters_for_mechanism} & \small Parameters in those equations. \\
\quad\fn{mechanisms_for_observation} & \small Mechanisms linked to an observation entity (schema-ready). \\
\quad\fn{documents_supporting} & \small Documents that support an entity. \\
\quad\fn{sources_for_relationship} & \small Provenance of an edge. \\
\quad\fn{evidence_connecting} & \small Relationships and shared equations linking two variables. \\
\addlinespace[4pt]
\multicolumn{2}{@{}l}{\small\bfseries\texttt{knowledge\_graph/schema.py}}\\[1pt]
\quad\fn{EntityType} & \small Allowed entity types (Document, Mechanism, Variable, Parameter, Equation, Observation). \\
\quad\fn{RelationType} & \small Allowed relationship types: GOVERNS, HAS\_VARIABLE, HAS\_PARAMETER, SUPPORTED\_BY, CAUSES, DEPENDS\_ON, and others. \\
\quad\fn{Entity} & \small A node: id, type, name, metadata. \\
\quad\fn{SourceRef} & \small Where a fact came from: document and chunk. \\
\quad\fn{Relationship} & \small A typed edge with its sources. \\
\addlinespace[4pt]
\multicolumn{2}{@{}l}{\small\bfseries\texttt{knowledge\_graph/serialization.py}}\\[1pt]
\quad\fn{graph_to_dict} & \small Convert the graph to a JSON-able dict. \\
\quad\fn{graph_from_dict} & \small Rebuild a graph from that dict. \\
\quad\fn{save_graph} & \small Write the graph to a JSON file. \\
\quad\fn{load_graph} & \small Read a graph from a JSON file. \\
\addlinespace[4pt]
\end{longtable}
\subsection*{\texttt{evaluation/}}
\begin{longtable}{@{}>{\raggedright\arraybackslash}p{5.6cm}>{\raggedright\arraybackslash}p{10.6cm}@{}}
\multicolumn{2}{@{}l}{\small\bfseries\texttt{evaluation/ablation.py}}\\[1pt]
\quad\fn{_discovery_only} & \small The SINDy-alone arm: discover, then match the result to a documented mechanism. \\
\quad\fn{run_ablation} & \small Run all seven arms on the 12 datasets. \\
\quad\fn{_median} & \small Median helper. \\
\quad\fn{summarize_ablation} & \small Per-arm totals: correct model, wrong candidates rejected, parameter and OOD accuracy. \\
\addlinespace[4pt]
\multicolumn{2}{@{}l}{\small\bfseries\texttt{evaluation/adversarial.py}}\\[1pt]
\quad\fn{_CannedProvider} & \small Test LLM that returns a fixed (possibly malicious) response. \\
\quad\fn{_CannedProvider.generate} & \small Return the canned response. \\
\quad\fn{_H} & \small Shorthand to build a test Hypothesis. \\
\quad\fn{_case} & \small Record one adversarial case result. \\
\quad\fn{find_leaks} & \small Search corpus text for any hidden benchmark value. \\
\quad\fn{run_adversarial} & \small Run the 20 attack cases. \\
\addlinespace[4pt]
\multicolumn{2}{@{}l}{\small\bfseries\texttt{evaluation/benchmark.py}}\\[1pt]
\quad\fn{_item} & \small Record one question result. \\
\quad\fn{_hyp} & \small Build a Hypothesis from an equation string, inferring parameter names. \\
\quad\fn{run_single_questions} & \small The 50 single questions over the 12 datasets and the graph. \\
\quad\fn{run_multihop} & \small The 25 multi-hop chains. \\
\addlinespace[4pt]
\multicolumn{2}{@{}l}{\small\bfseries\texttt{evaluation/datasets.py}}\\[1pt]
\quad\fn{Dataset} & \small One benchmark system: family, truth, grids, field, initial condition, exact solution. \\
\quad\fn{_diffusion} & \small Exact two-mode diffusion dataset. \\
\quad\fn{_gauss} & \small Exact moving, spreading Gaussian; used for advection (D = 0) and advection-diffusion. \\
\quad\fn{_advection} & \small Pure advection dataset. \\
\quad\fn{_adv_diff} & \small Advection-diffusion dataset. \\
\quad\fn{_reaction_diffusion} & \small Reaction-diffusion dataset solved numerically with RK4. \\
\quad\fn{benchmark_datasets} & \small The 12 datasets (4 families x 3 settings), cached to .npz. \\
\quad\fn{_build_all} & \small Construct all 12 datasets. \\
\quad\fn{_rd_shell} & \small Reaction-diffusion dataset without its (expensive) field, for quick metadata. \\
\quad\fn{add_noise} & \small Add Gaussian noise to a dataset's field. \\
\addlinespace[4pt]
\multicolumn{2}{@{}l}{\small\bfseries\texttt{evaluation/describe.py}}\\[1pt]
\quad\fn{_centroid} & \small Amplitude-weighted centre of a profile. \\
\quad\fn{_spread} & \small Width of a profile. \\
\quad\fn{_n_peaks} & \small Number of peaks in a profile. \\
\quad\fn{observation_stats} & \small Amplitude change, drift, spread and peaks between early and late times. \\
\quad\fn{describe_observations} & \small Fixed-template text description from those statistics (never names a mechanism). \\
\addlinespace[4pt]
\multicolumn{2}{@{}l}{\small\bfseries\texttt{evaluation/engine.py}}\\[1pt]
\quad\fn{Resources} & \small Corpus, RAG index and graph, built once and shared. \\
\quad\fn{build_pool} & \small Description -> RAG/graph context -> candidate pool (evidence equations, graph templates, mock LLM). \\
\quad\fn{_grid_param_contributions} & \small Term necessity computed on grid derivatives instead of a PINN. \\
\quad\fn{_reduced_term_set} & \small Term set of an equation with its inactive terms removed. \\
\quad\fn{validate_and_select} & \small Fast pipeline: least-squares fit per candidate, gate, selection. \\
\quad\fn{predict_ood} & \small Solve the selected law forward beyond the data window. \\
\quad\fn{family_of} & \small Map a hypothesis's term set to its mechanism family. \\
\quad\fn{truth_field} & \small Ground-truth field at given times (evaluation only). \\
\quad\fn{score_run} & \small Score one run: correct family, parameter error, OOD error. \\
\quad\fn{family_of_row} & \small Mechanism family from a stored result row. \\
\addlinespace[4pt]
\multicolumn{2}{@{}l}{\small\bfseries\texttt{evaluation/identify.py}}\\[1pt]
\quad\fn{_syms} & \small SymPy symbols for the fields. \\
\quad\fn{term_set} & \small The set of field terms in a hypothesis. \\
\quad\fn{_label_to_expr} & \small Convert a library label like 'u*u\_x' to a SymPy expression. \\
\quad\fn{discovered_coefficients} & \small Non-zero SINDy coefficients as \{term: value\}. \\
\quad\fn{match_mechanism} & \small Find the documented template with exactly the discovered term set. \\
\quad\fn{recover_parameters} & \small Solve for a template's parameters from discovered coefficients. \\
\quad\fn{grid_fields} & \small Finite-difference fields of a gridded dataset, flattened. \\
\quad\fn{fit_candidate_on_grid} & \small Least-squares fit of u\_t on a candidate's own terms; train and holdout error. \\
\quad\fn{library_collinearity} & \small Identifiability check: maximum correlation between library columns. \\
\addlinespace[4pt]
\multicolumn{2}{@{}l}{\small\bfseries\texttt{evaluation/leakage.py}}\\[1pt]
\quad\fn{ground_truth_references} & \small List ground-truth identifiers mentioned in a function's source code. \\
\quad\fn{decision_functions} & \small The six functions that make decisions (pool, validation, splits, selection). \\
\quad\fn{leakage_audit} & \small Static audit of those functions plus a corpus scan for hidden values. \\
\addlinespace[4pt]
\end{longtable}
\subsection*{\texttt{experiments/}}
\begin{longtable}{@{}>{\raggedright\arraybackslash}p{5.6cm}>{\raggedright\arraybackslash}p{10.6cm}@{}}
\multicolumn{2}{@{}l}{\small\bfseries\texttt{experiments/02\_equation\_discovery.py}}\\[1pt]
\quad\fn{load_config} & \small Read config.yaml. \\
\quad\fn{run_case} & \small Run discovery on one dataset variant and evaluate it. \\
\quad\fn{main} & \small Phase 2: clean, sparse, noisy (raw and smoothed); save report, table and plot. \\
\addlinespace[4pt]
\multicolumn{2}{@{}l}{\small\bfseries\texttt{experiments/03\_forward\_pinn.py}}\\[1pt]
\quad\fn{load_config} & \small Read config.yaml. \\
\quad\fn{main} & \small Phase 3 forward PINN with D = 0.1 fixed. \\
\addlinespace[4pt]
\multicolumn{2}{@{}l}{\small\bfseries\texttt{experiments/04\_inverse\_pinn.py}}\\[1pt]
\quad\fn{load_config} & \small Read config.yaml. \\
\quad\fn{main} & \small Phase 3 inverse PINN estimating D from 0.2. \\
\addlinespace[4pt]
\multicolumn{2}{@{}l}{\small\bfseries\texttt{experiments/05\_hypothesis\_validation.py}}\\[1pt]
\quad\fn{load_config} & \small Read config.yaml. \\
\quad\fn{run_hypothesis} & \small Train one hypothesis's PINN and score it. \\
\quad\fn{main} & \small Phase 3: diffusion vs advection with a time-based OOD split. \\
\addlinespace[4pt]
\multicolumn{2}{@{}l}{\small\bfseries\texttt{experiments/06\_hypothesis\_generation.py}}\\[1pt]
\quad\fn{load_config} & \small Read config.yaml. \\
\quad\fn{main} & \small Phase 4: generate, validate and rank hypotheses; optionally run PINNs. \\
\addlinespace[4pt]
\multicolumn{2}{@{}l}{\small\bfseries\texttt{experiments/07\_scientific\_rag\_kg.py}}\\[1pt]
\quad\fn{load_config} & \small Read config.yaml. \\
\quad\fn{hyp_to_dict} & \small Serialise a hypothesis. \\
\quad\fn{summarize_generation} & \small Counts for a generation run (candidates, valid, duplicates, evidence-supported). \\
\quad\fn{main} & \small Phase 5: corpus, RAG, graph, context, leakage check, Mode A vs Mode B. \\
\addlinespace[4pt]
\multicolumn{2}{@{}l}{\small\bfseries\texttt{experiments/08\_physics\_constrained\_search.py}}\\[1pt]
\quad\fn{load_config} & \small Read config.yaml. \\
\quad\fn{stage_prepare} & \small Phase 6 stage 1: context, generation, pool, filters. \\
\quad\fn{_load_survivors} & \small Reload the filtered candidates. \\
\quad\fn{_splits} & \small Build train/holdout/OOD splits from noisy.npz. \\
\quad\fn{_setup} & \small Domain and initial-condition settings for the PINN. \\
\quad\fn{stage_pinn} & \small Phase 6 stage 2: inverse PINN per survivor, cached. \\
\quad\fn{oracle_evaluation} & \small Benchmark-only error against the hidden truth, after selection. \\
\quad\fn{stage_select} & \small Phase 6 stage 3: stability, selection, baselines, OOD, novel prediction, report. \\
\quad\fn{_plots} & \small Draw the Phase 6 summary figure. \\
\quad\fn{_report} & \small Write phase6\_report.md. \\
\quad\fn{main} & \small Run the chosen stages. \\
\addlinespace[4pt]
\multicolumn{2}{@{}l}{\small\bfseries\texttt{experiments/09\_evaluation.py}}\\[1pt]
\quad\fn{save} & \small Write one Phase 7 result file. \\
\quad\fn{real_llm} & \small Mode A vs Mode B with a real LLM on 3 datasets (only with -{}-real-llm and a key). \\
\quad\fn{acc} & \small Format an accuracy fraction. \\
\quad\fn{report} & \small Write phase7\_report.md. \\
\quad\fn{main} & \small Phase 7: questions, multi-hop, adversarial, ablation, leakage, determinism check, report. \\
\addlinespace[4pt]
\multicolumn{2}{@{}l}{\small\bfseries\texttt{experiments/10\_generalization.py}}\\[1pt]
\quad\fn{_truth_dataset} & \small The hidden advection-diffusion system (c = 0.4, D = 0.02). \\
\quad\fn{observations} & \small 400 noisy samples, the only data the pipeline sees. \\
\quad\fn{_recon} & \small Grid reconstruction for SINDy and the description. \\
\quad\fn{_setup} & \small Domain settings for the PINN. \\
\quad\fn{stage_prepare} & \small SINDy baseline, pool and filters for the new system. \\
\quad\fn{_survivors} & \small Reload filtered candidates. \\
\quad\fn{_splits} & \small Train/holdout/OOD splits. \\
\quad\fn{stage_pinn} & \small Inverse PINN per candidate. \\
\quad\fn{stage_select} & \small Selection under both rules, OOD, oracle evaluation. \\
\quad\fn{main} & \small Run the chosen stages. \\
\addlinespace[4pt]
\multicolumn{2}{@{}l}{\small\bfseries\texttt{experiments/11\_real\_llm\_gemini.py}}\\[1pt]
\quad\fn{_load_script} & \small Import an experiment script by file path. \\
\quad\fn{_phase6_module} & \small Load the Phase 6 script and redirect its outputs to the Gemini folder. \\
\quad\fn{stage_generate} & \small Run 09's real\_llm Mode A vs Mode B with Gemini. \\
\quad\fn{annotate_report} & \small Mark the reused Phase 6 report as a real-LLM run. \\
\quad\fn{compare_with_mock} & \small Side-by-side summary of the mock and Gemini Phase 6 runs. \\
\quad\fn{main} & \small Check the key, force the Gemini provider, run the chosen stages. \\
\addlinespace[4pt]
\multicolumn{2}{@{}l}{\small\bfseries\texttt{experiments/12\_conservation.py}}\\[1pt]
\quad\fn{_windows} & \small Conservation violation in the train, OOD and beyond-data windows. \\
\quad\fn{evaluate_run} & \small Conservation check for every saved candidate of one Phase 6 run (no retraining). \\
\quad\fn{main} & \small Check the mock and Gemini runs plus the physics-free baseline; write results/evaluation/conservation/. \\
\addlinespace[4pt]
\end{longtable}
